\label{ch:conclusiones}

El desarrollo de este proyecto permitió diseñar, fabricar y validar una 
impresora 3D tipo FDM de arquitectura abierta, cumpliendo satisfactoriamente 
con el objetivo general planteado. A partir del análisis de los resultados 
experimentales, se presentan las siguientes conclusiones:

\begin{itemize}
    \item \textbf{Diseño y fabricación:} Se construyó un equipo funcional y 
    modular, utilizando componentes de código abierto y manufactura aditiva. 
    La integración de una estructura de perfiles de aluminio y el uso de 
    firmware Marlin permitieron obtener una plataforma robusta, cuya 
    calibración y mantenimiento están al alcance de un estudiante.

    \item \textbf{Rediseño del sistema de extrusión y enfriamiento:} La 
    sustitución de los componentes originales por un sistema de accionamiento 
    directo (\textit{direct drive}) con una base estructural de aluminio 
    resolvió los problemas de incompatibilidad térmica. Se logró una reducción 
    del 24.2\% en la masa del módulo, lo que disminuyó la inercia sobre el 
    eje \(X\). Asimismo, la incorporación de la base de aluminio y el 
    ventilador redujo la temperatura del disipador en un 42.25\%, mitigando 
    el riesgo de \textit{heat creep} y garantizando la continuidad del flujo 
    de filamento.

    \item \textbf{Calidad y precisión dimensional:} El sistema demostró ser 
    capaz de fabricar piezas con errores dimensionales inferiores al 0.3\%, 
    situándose dentro de los estándares de tolerancia para impresoras FDM de 
    escritorio. La implementación de mejoras estructurales, como las escuadras 
    de refuerzo y los tensores ajustables, permitió mitigar defectos como el 
    \textit{ghosting}.

    \item \textbf{Diagnóstico de fallas y materiales:} Se comprobó que el 
    sistema es altamente sensible a las condiciones del filamento. A través 
    de un proceso de diagnóstico experimental, se identificó la humedad 
    ambiental como la causa de la fragilidad y el acabado irregular en las 
    piezas, validando un método de secado por cama calefactada como una 
    solución provisional efectiva. Asimismo, la torre de temperatura permitió 
    establecer que la ventana térmica óptima para el PLA utilizado se sitúa 
    en \qty{205}{\celsius}.

    \item \textbf{Viabilidad económica y didáctica:} El costo total de 
    fabricación (\$4189 MXN) resultó menor que el de impresoras comerciales 
    de gama baja. Si bien el equipo carece de automatizaciones (como 
    nivelación automática), esta característica se convierte en una ventaja 
    pedagógica, ya que obliga al usuario a comprender y ajustar manualmente 
    cada subsistema. La modularidad y capacidad de autorreparación confirman 
    la vigencia de la filosofía RepRap en entornos académicos con recursos 
    limitados.
\end{itemize}

En términos generales, el proyecto demuestra que es posible construir una 
herramienta de manufactura digital funcional y precisa mediante la 
investigación aplicada, el diseño iterativo y el aprovechamiento de 
tecnologías abiertas. Las limitaciones identificadas, como la velocidad de 
impresión o la falta de ciertos sensores, no demeritan el desempeño del 
equipo, sino que establecen una hoja de ruta clara para futuras generaciones 
de estudiantes que deseen continuar con el desarrollo de esta plataforma.